% !TeX encoding = UTF-8
% !TeX root = ../qjbc-talk.tex
% 第二项工作选题汇报：10 个内容帧（封面与结尾由入口提供）
% 口径：组会说明"为什么做这个 / 受了什么启发 / 可行性到什么程度"
% 不在本稿展开的内容：阈值裕量目标本身的实验臂结果（E009 的 W_SOFTMIN / W_MEAN_MARGIN
% 三臂对照）。该结果不适合作为选题依据写入汇报，故本稿只保留由 E007 内部对比得出的定性结论。

% ============================================================
% 1. 场景与一句话选题
% ============================================================
\begin{frame}{选题一句话：一个攻击，两个安全目标}
  \textbf{攻击者设定}：已登记的低权限用户，系统里保存着他的源身份模板；
  他要生成一张自然人脸，同时做到
  \vspace{0.3em}
  \begin{itemize}\setlength{\itemsep}{2pt}
    \item \textbf{Privilege escalation（越权冒充）}：认证入口接受他为\textbf{高权限目标身份}
    \item \textbf{Attribution evasion（归因规避）}：审计入口\textbf{不再关联回他自己的源身份}
  \end{itemize}
  \vspace{0.3em}
  给定目标 FR 模型 $f_v$ 与阈值 $\tau_v$，把相似度换算成\textbf{裕量}（分数离及格线的距离，
  正=过线、负=没过）：
  \[
  m_t = s_t-\tau_v \;\text{（target accept）},
  \qquad
  m_s = \tau_v-s_s \;\text{（source reject）}
  \]
  \[
  \boxed{\,m_{\mathrm{joint}}=\min(m_t,m_s)>0
  \;\Longleftrightarrow\;
  \text{联合攻击成功}\,}
  \]
  {\small 这不是单目标 impersonation：要求两个条件\textbf{同时跨过各自的绝对阈值}，
  而阈值由\textbf{别人的模型}决定，攻击者看不到。\textbf{选题}：用本地扩散模型生成自然候选，
  再用\textbf{极少量}黑盒查询校正代理边界与真实边界之间的偏差。}
\end{frame}

% ============================================================
% 2. 为什么做（一）：单目标 ASR 会把联合失败藏起来
% ============================================================
\begin{frame}{为什么做（一）：单目标 ASR 会把联合失败藏起来}
  \vspace{-4pt}
  \footnotesize
  \begin{columns}[T]
    \column{.49\textwidth}
    \begin{block}{工作一（方向性铰链项）}
      {\footnotesize
      \textbf{现象}：FaceNet 上 ASR 77\% 时，真正身份替换仅 38\%（工作一，400 对，FAR=0.01）。
      多数样本\textbf{同时挂着两个身份}——target 接受了，源身份没甩掉。\\[0.3em]
      \textbf{机制}：只要求 target 比 source 高出 $m$，高出后惩罚归零，
      \textbf{不要求 source 落到阈值以下}：\par\vspace{0.3em}
      $\mathcal L_{1}^{(k)}=\lambda_t(1-S_{\mathrm{tgt}}^{(k)})
      +\lambda_s[S_{\mathrm{src}}^{(k)}+m-S_{\mathrm{tgt}}^{(k)}]_+$}
    \end{block}

    \column{.49\textwidth}
    \begin{block}{DiffAIM（差值打分）}
      {\footnotesize
      \textbf{同构问题}：目标函数是每代理的总分\par\vspace{0.2em}
      $\mathrm{score}_k=S_{\mathrm{tgt}}^{(k)}-S_{\mathrm{src}}^{(k)}+1$\par\vspace{0.3em}
      含 $-S_{\mathrm{src}}$ 压源项，但\textbf{只优化相对差值}：只要 target 比 source 高，
      分数就好看，\textbf{两个分数的绝对位置无人管}。}
    \end{block}
  \end{columns}
  \vspace{0.2em}

  \begin{beamercolorbox}[wd=\textwidth,sep=3pt]{block body alerted}
    \footnotesize\textbf{共同缺口：} 两个目标都是"排序 / 差值"结构，都不含\textbf{决策阈值}。
    ASR 只统计 target 侧是否过线，因此"target 过、source 也过"的样本照样计入成功。
    真正决定攻击成败的是 $\max\min(m_t,m_s)$，而不是平均分数。
  \end{beamercolorbox}
\end{frame}

% ============================================================
% 3. 为什么做（二）：同一批脸换评测模型，失败形态整体换了一种
% ============================================================
\begin{frame}{为什么做（二）：换一个评测模型，失败形态就换了}
  \vspace{-6pt}
  \footnotesize
  DiffAIM 在同一批 \textbf{1000 对} CelebA 样本上，两个未参与训练的黑盒模型按各自阈值
  分类（FAR=0.01：FaceNet $\tau=0.409$、AdaFace $\tau=0.145$）：
  \vspace{0.2em}

  \centering
  \setlength{\tabcolsep}{6pt}\renewcommand{\arraystretch}{1.05}
  \begin{tabular}{lcccc}
    \toprule
    黑盒模型 & 联合成功（tgt过、src未过） & 双过（身份叠加） & 仅 src 过 & 双不过 \\
    \midrule
    FaceNet & \textbf{612} & \textbf{37} & 44 & 307 \\
    AdaFace & \textbf{605} & \textbf{267} & 66 & 62 \\
    \bottomrule
  \end{tabular}

  \vspace{0.3em}
  \raggedright
  {\footnotesize
    \textbullet\ 联合成功数几乎一样（612 vs 605）：\textbf{代理侧学到的方向是有效的}；\newline
    \textbullet\ 但叠加数从 \textbf{37 涨到 267（约 7 倍）}，两个模型同时判叠加的只有 \textbf{26 对}；\newline
    \textbullet\ 同一张脸，换个模型就从"恰好过线"变成"两个身份都匹配"。
  }

  \vspace{0.35em}
  \begin{beamercolorbox}[wd=\textwidth,sep=3pt]{block body example}
    \centering\footnotesize
    \textbf{结论（这是选题的直接理由）：} 相对差值目标不约束分数的绝对位置，
    所以判据一换模型，边界就漂。缺的不是"更强的扰动"，是\textbf{阈值/边界这一层}——
    攻击必须在\textbf{真实 victim 的两个绝对阈值}上同时成立。
  \end{beamercolorbox}
\end{frame}

% ============================================================
% 4. 选题怎么来的：一条失败链，不是一次灵感
% ============================================================
\begin{frame}{选题怎么来的：一条失败链，不是一次灵感}
  \vspace{-8pt}
  \scriptsize
  \renewcommand{\arraystretch}{1.1}
  \setlength{\tabcolsep}{3pt}
  \begin{tabular}{p{0.7cm}p{4.0cm}p{4.7cm}p{3.6cm}}
    \toprule
    编号 & 当时想做的事 & 结果 & 对下一步的约束 \\
    \midrule
    E004 & 联合选择代理子集与等面积支撑域提升迁移 & 区域效应存在，但代理错配与可利用结构未过 Gate B/C $\to$ stop & 仅靠“选配置”解决不了迁移 \\
    E005 & FR 判定与真人身份感知是否有可利用错位 & 标注质量门未过，正式结论不可评价 $\to$ stop & 不再押在高成本人类测量上 \\
    E006 & 冒充与逃逸是否共享可路由的控制支撑 & 局部连续分数有信号，端到端被直接联合优化支配，joint 成功 0 $\to$ stop & “改哪里”不是瓶颈，“优化什么”才是 \\
    E007 & 把平均双 loss 换成\textbf{最差联合裕量} & 代理侧明显改善、留出同类模型也改善，但 FaceNet 仍 0 跨界 & 代理可行 $\neq$ 目标跨界，剩余是边界偏差 \\
    E008 & 只开后期两个低维变量，用少量 victim 查询校正 & 跨界在小样本上出现（见可行性页） & 值得进入方法开发，但主张必须先收窄 \\
    \bottomrule
  \end{tabular}

  \vspace{0.3em}
  \begin{beamercolorbox}[wd=\textwidth,sep=3pt]{block body alerted}
    \footnotesize
    这条链的逻辑不是"某方法有效所以放大它"，而是\textbf{每走一步就排除一个解释}：
    不是区域没选对，不是人类感知错位，不是"改哪里"不够，也不只是优化没收敛。
    剩下的唯一解释是\textbf{代理边界与真实边界不重合}。
  \end{beamercolorbox}
\end{frame}

% ============================================================
% 5. 关键诊断：代理侧可行，目标侧不跨界
% ============================================================
\begin{frame}{关键诊断：代理侧已经修好，victim 仍不跨界}
  \vspace{-8pt}
  \scriptsize
  \textbf{E007 设定}：15 个身份对（5 个 target 身份簇）× 3 个 LOSO fold × 3 个目标函数 = 135 个 arm，
  每个 arm 固定 300 次 backward，\textbf{victim 查询数 0}。J=平均铰链、WM=裕量求平均、
  W=裕量 soft-min（最差约束）。

  \vspace{0.3em}
  \centering
  \setlength{\tabcolsep}{5pt}\renewcommand{\arraystretch}{1.05}
  \begin{tabular}{lccc}
    \toprule
    方法 & 训练侧最差联合裕量中位 & 训练侧 joint 成功 & FaceNet joint 成功 \\
    \midrule
    J（平均铰链） & $-0.0192$ & 19/45 & 0/45 \\
    WM（裕量平均） & $-0.0926$ & 11/45 & 0/45 \\
    \textbf{W（soft-min 最差裕量）} & $\mathbf{+0.0694}$ & \textbf{36/45} & \textbf{0/45} \\
    \bottomrule
  \end{tabular}

  \vspace{0.35em}
  \raggedright
  {\footnotesize
    \textbullet\ 训练侧 36/45 $\to$ FaceNet 0/45，这是清楚的 \textbf{train--victim gap}；\newline
    \textbullet\ 留出同类代理（LOSO）的 joint 裕量在 \textbf{15/15 个身份对}上都比 J 改善
    （中位差 $+0.0606$，簇 bootstrap 95\% CI $[+0.039,+0.087]$）；\newline
    \textbullet\ 但留出侧绝对成功率只有 1/45，FaceNet joint 裕量中位仍 $-0.2406$，
    且 W$-$J 的 FaceNet joint 裕量 CI \textbf{跨 0}；\newline
    \textbullet\ 所以不能说"W 已经解决黑盒迁移"。\textbf{它只证明代理侧可行域可以修好、还能外推。}
  }

  \vspace{0.3em}
  \begin{beamercolorbox}[wd=\textwidth,sep=3pt]{block body example}
    \centering\footnotesize
    问题定位从"优化不收敛"变成"\textbf{代理可行域与真实决策边界不重合}"——
    这正是后续方法的入口。
  \end{beamercolorbox}
\end{frame}

% ============================================================
% 6. 剩余缺口有多大：失败样本贴着边界
% ============================================================
\begin{frame}{剩余缺口有多大：一半以上失败样本贴着边界}
  \vspace{-6pt}
  \scriptsize
  \textbf{为什么不继续加强代理侧裕量优化}：E007 内部对比 W vs WM，soft-min 越强，修补越集中于
  source-reject 一侧——留出侧与 FaceNet 上都有 \textbf{36/45} 个 fold 的 target 裕量反而下降。
  继续加步数只会加剧单侧偏置，不会产生跨界。
  \vspace{0.25em}

  \textbf{DiffAIM 在 FaceNet 上的 1000 对：失败样本的联合裕量分布}
  \vspace{0.2em}

  \centering
  \setlength{\tabcolsep}{6pt}\renewcommand{\arraystretch}{1.05}
  \begin{tabular}{lcccc}
    \toprule
    联合裕量区间 & $[-0.02,0)$ & $[-0.05,0)$ & $[-0.10,0)$ & 全部失败 \\
    \midrule
    样本数 & \textbf{38} & \textbf{107} & \textbf{209} & 388 \\
    占失败比例 & 9.8\% & 27.6\% & \textbf{53.9\%} & 100\% \\
    覆盖 target 簇 & 5/5 & 5/5 & 5/5 & — \\
    \bottomrule
  \end{tabular}

  \vspace{0.3em}
  \raggedright
  {\footnotesize
    \textbullet\ 一半以上失败样本离边界\textbf{不到 0.10}，而且分布在全部 5 个 target 簇上，
    不是某一类的特例；\newline
    \textbullet\ 这些样本需要的不是"重做一遍优化"，而是\textbf{最后一小段位移}；\newline
    \textbullet\ 但反过来说：能补的是\textbf{近边界群体}，不是任意一张脸。
  }

  \vspace{0.3em}
  \begin{beamercolorbox}[wd=\textwidth,sep=3pt]{block body example}
    \centering\footnotesize
    \textbf{缺口位置清楚了}：既然只差一点，且差的方向只有 victim 自己知道，
    那就可以用\textbf{少量真实反馈}去补，而不是继续猜。
  \end{beamercolorbox}
\end{frame}

% ============================================================
% 7. 受什么启发：四条成熟路线，缺的是接到双阈值判据上
% ============================================================
\begin{frame}{受了什么启发：四条成熟路线，缺的是接线}
  \vspace{-8pt}
  \scriptsize
  \renewcommand{\arraystretch}{1.42}
  \setlength{\tabcolsep}{3pt}
  \begin{tabular}{p{2.9cm}p{3.05cm}p{4.1cm}p{2.95cm}}
    \toprule
    技术路线 & 代表工作 & 已证明的事 & 本工作怎么用 \\
    \midrule
    迁移先验+目标反馈 & TREMBA、LeBA &
      代理搜索空间 + victim 反馈，比从零查询省 &
      代理生成候选，反馈只修偏差\\[3pt]
    低维结构化搜索 & AutoZOOM、DifAttack、QEBA、GADA &
      查询空间应按任务结构降维 &
      只动后期两个引导系数\\[3pt]
    FR 边界搜索 & Evol.~Attack、RAF &
      FR 边界可在低维语义参数上搜 &
      语义参数换成轨迹参数\\[3pt]
    少量变量在线调节 & Wei et al.-TPAMI'23 &
      少量变量 + 好初始化就能压到十级查询 &
      只开 $b$、$\eta$ 两个变量\\[3pt]
    目标反馈替代代理 & PPOM-Attack &
      直接用 victim 反馈避开代理偏差 &
      动机成立（其预算千级）\\[3pt]
    扩散生成的自然性 & DiffAttack、DiffAM、CLIP2Protect &
      生成流形本身就是质量先验 &
      冻结反演，只重放后缀 \\
    \bottomrule
  \end{tabular}

  \vspace{0.3em}
  \begin{beamercolorbox}[wd=\textwidth,sep=2pt]{block body alerted}
    \footnotesize
    \textbf{诚实的定位}：这些都是成熟范式，"低维 + 查询校正"本身不是新东西。本地语料
    没有直接做"扩散轨迹上双身份联合阈值校正"的论文，但正式新颖性审计尚未完成，
    汇报中\textbf{不使用"首次提出"}；要回答的是\textbf{这套范式能否接到双阈值判据上}。
  \end{beamercolorbox}
\end{frame}

% ============================================================
% 8. 方法：冻结轨迹，只动两个数
% ============================================================
\begin{frame}{拟议方法：冻结轨迹，只动两个数}
  \vspace{-6pt}
  \footnotesize
  \textbf{设计前提（由前五页逼出）}：代理侧可修但不可靠、失败贴着边界、裕量不能当进攻目标
  $\Rightarrow$ 跨界只能靠\textbf{目标侧反馈}完成，且反馈必须落在\textbf{少量、低维}的参数上。
  \vspace{0.3em}

  \textbf{载体：冻结反演，只重放扩散后期}
  \begin{itemize}\setlength{\itemsep}{1pt}\footnotesize
    \item 复用 DiffAIM 的 DDIM 反演（100 步、skip=80 $\to$ 20 个引导时间步）
    \item 前缀 $t{=}191\!\to\!61$ 固定（240 次反向），只重放后缀 $t{=}51\!\to\!1$（60 次）
    \item 图像自然性由生成器保证，不是靠像素约束"抠"出来
  \end{itemize}
  \vspace{0.2em}

  \textbf{控制变量：只有两个}
  \begin{itemize}\setlength{\itemsep}{1pt}\footnotesize
    \item $\eta$：后期更新强度；$b$：target 侧 vs source 侧的角色权重偏置
    \item 9 点网格 $\times$ 2 个相似度（target + source）= \textbf{每样本 18 次查询}
    \item 每个候选只向 victim 取两个分数，直接算 $m_{\mathrm{joint}}$，不做梯度估计
  \end{itemize}
  \vspace{0.2em}

  \textbf{选择规则预先固定，不做事后挑选}
  \begin{itemize}\setlength{\itemsep}{1pt}\footnotesize
    \item 用 victim 的 $m_{\mathrm{joint}}$ 选点，同时报告"若只用代理信号会选到哪个点"
    \item 停止条件：跨界 / 用完预算 / 质量约束触发；成功与失败样本都报告
  \end{itemize}
\end{frame}

% ============================================================
% 9. 可行性（一）：18 次查询能推过一部分
% ============================================================
\begin{frame}{可行性（一）：18 次查询能推过一部分近边界样本}
  \vspace{-8pt}
  \scriptsize
  冻结反演与轨迹，只重放后期时间窗，每样本 18 次查询；样本取自 E007 面板并按 FaceNet
  联合裕量排序（\textbf{近边界群体}）。
  \vspace{0.25em}

  \centering
  \setlength{\tabcolsep}{5pt}\renewcommand{\arraystretch}{1.02}
  \begin{tabular}{llccc}
    \toprule
    组别 & 样本与网格 & 跨界 & 平均 $\Delta m_v$ & victim 选择器优于代理信号 \\
    \midrule
    E008 & 近边界开发集，$\eta\le1.5$ & 0/5 & $+0.0263$ & 4/5 \\
    E008b & 同上，$\eta$ 外扩到 3.0 & 1/5 & $+0.0321$ & 3/5 \\
    E008c & 远离边界（选样失误，不作为泛化证据） & 0/5 & $+0.0704$ & 5/5 \\
    \textbf{E008d} & \textbf{未参与选样的近边界集} & \textbf{2/5} & $\mathbf{+0.0792}$ & \textbf{5/5} \\
    \bottomrule
  \end{tabular}

  \vspace{0.3em}
  \raggedright
  {\footnotesize
    \textbullet\ \textbf{跨界可复现}：E008b 的 1/5 在未参与选样的 E008d 上为 2/5，累计 \textbf{3/20}；\newline
    \textbullet\ \textbf{幅度可观}：p00230/f2 从 $-0.1086$ 推到 $+0.0047$，单样本改善 $+0.1133$；\newline
    \textbullet\ \textbf{反馈不可替代}：E008d 上 victim 选择器 5/5 优于"只用代理信号"；\newline
    \textbullet\ \textbf{质量守住}：最优点 LPIPS 0.02–0.04，与未校正时同区间；\newline
    \textbullet\ \textbf{诚实标注}：$\eta=3.0$ 上界被 5/5 选为最优，网格尚未收敛；
    跨界率是\textbf{近边界群体}的小样本数字，不是总体 ASR。
  }

  \vspace{0.2em}
  \begin{beamercolorbox}[wd=\textwidth,sep=3pt]{block body example}
    \centering\footnotesize
    \textbf{每样本 18 次查询，可以把一部分原本必败的近边界样本推过判定边界。}\\
    {\scriptsize 不构成与 released DiffAIM / J 的同预算方法比较。}
  \end{beamercolorbox}
\end{frame}

% ============================================================
% 10. 可行性（二）：工程底座、风险与下一步
% ============================================================
\begin{frame}{可行性（二）：底座已就位，风险已列清}
  \vspace{-8pt}
  \scriptsize
  \begin{columns}[T]
    \column{.47\textwidth}
    {\scriptsize\textbf{已具备的工程基础}
    \begin{itemize}\setlength{\itemsep}{0pt}
      \item 扩散反演、生成、代理集成、质量度量已跑通
      \item J / WM / W 三种目标与 LOSO runner 已实现
      \item manifest、resume、聚合、trace 已完成（135 arm 验证过）
      \item 单 arm 约 49 秒；查询账本可逐条审计
      \item 在线阶段只更新 2 个标量，不重训生成器
    \end{itemize}}

    \column{.51\textwidth}
    {\scriptsize\textbf{主要风险与控制}
    \begin{itemize}\setlength{\itemsep}{0pt}
      \item 审计侧可能是硬标签 $\to$ 明确 score-based 假设，decision-only 列为扩展
      \item 低维空间表达不足 $\to$ 与更宽时间窗 / 随机同维子空间做 oracle 诊断
      \item 通用黑盒优化即可达成 $\to$ 同维度同预算 SPSA/NES/CMA-ES 正面对比
      \item 起点无必要（代理已饱和）$\to$ 三种初始化进同一校正器
      \item victim 过拟合 / 画质换跨界 $\to$ 调参与测试 victim 分离，固定 LPIPS 预算
    \end{itemize}}
  \end{columns}

  \vspace{0.2em}
  {\scriptsize\textbf{下一步（按顺序）}
  \begin{itemize}\setlength{\itemsep}{0pt}
    \item \textbf{查询数扫描}：读现有查询账本，画 $k{=}3/6/9/12/18$ 累计跨界曲线（不占算力）
    \item \textbf{外扩网格}：$\eta$ 到 4.0–4.5；$b$ 朝\textbf{正}方向扩（此前只朝负扩）
    \item \textbf{扩样本}：5 例扩到 15 例，覆盖全部 5 个 target 簇；\textbf{正面对照}：同 18 次预算下 DiffAIM 起点 vs W 起点
  \end{itemize}}

  \vspace{0.15em}
  \begin{beamercolorbox}[wd=\textwidth,sep=2pt]{block body example}
    \centering\footnotesize
    中心问题：\textbf{相同查询与质量预算下，少量 victim 反馈能否把“代理已改善、FaceNet 未跨界”
    的样本推过真实联合边界？}\par\vspace{0.15em}
    {\scriptsize 判定只看两项：是否产生新的 victim joint success；每次成功花了多少查询。}
  \end{beamercolorbox}
\end{frame}
